The traditional view of a neuron as a feature detector or an efficient encoder has difficulties in explaining the function of motor neurons and experimentally observed variable and context-dependent response properties of neurons. We put forward an alternative perspective, modeling each neuron as a feedback controller within a closed loop that includes other neurons and the external environment. Based on the recently developed Direct Data-Driven Control (DD-DC) approach, we propose a biologically plausible controller which implicitly identifies the dynamics of the rest of the loop, infers its latent state and optimizes control. The DD-DC model of a neuron accounts for multiple neurophysiological observations, including the switch from potentiation to depression in Spike-Timing-Dependent Plasticity (STDP) and its asymmetry; temporally extended feedforward and feedback neuronal filters and their adaptation to input statistics; imprecision of the neuronal spike-generation mechanism under constant input; as well as the prevalence of variability and/or noise in brain operation. The DD-DC neuron offers an alternative to the feedforward, instantaneously responding McCulloch-Pitts-Rosenblatt unit as a primitive for constructing biologically-inspired neural networks.